\section{Brauner Staub}
\stepcounter{ctrpoeminyear}
%\addcontentsline{toc}{section}{\thechapter.\arabic{ctrpoeminyear} \; Brauner Staub}

\begin{strophe}
\vers{Brauner Staub, Du qualvoll Elend,}
\vers{Du, der genießt die Warnung}
\vers{unzähliger Seelen.}
\end{strophe}

\begin{strophe}
\vers{Flüster mir zu, wer gab Dir den Ruf?}
\vers{Den Ruf eines Grauens, einer wüsten Tortur?}
\vers{Ein Ende ist Dir nicht nachgesagt.}
\vers{Denn einst seinen Anfang genommen}
\vers{soll es dich heimsuchen,}
\vers{komme Nacht, komme Tag.}
\end{strophe}

\begin{strophe}
\vers{Brauner Staub, Du qualvoll Elend,}
\vers{Du, der genießt die Neugier}
\vers{unzähliger Seelen.}
\end{strophe}

\begin{strophe}
\vers{Flüster mir zu, wer gab Dir den Ruf?}
\vers{Den Ruf eines Mysteriums, einer heimlichen Figur?}
\vers{Was macht Dich so mystisch?}
\vers{Wer verlieh Dir deine phantasievolle Natur?}
\vers{Einmal empfangen soll es dich wiegen}
\vers{in ein Land der Träume}
\vers{durch seinen liebevollen Schwur.}
\end{strophe}

\begin{strophe}
\vers{Brauner Staub, du qualvoll Elend,}
\vers{Du, der genießt die Ahnung}
\vers{Unzähliger Seelen.}
\end{strophe}

\phantom{a}

\phantom{a}

\phantom{a}

\phantom{a}

\phantom{a}

\phantom{a}

\begin{strophe}
\vers{Die, die Dich kennen, stehen mit Liebe und Hass,}
\vers{zu deinem Gewand, das glitzert und matscht.}
\vers{Sie preisen dein Tun}
\vers{und verzeihen nicht dein Getue.}
\vers{Für einen Moment des Glücks zerstören sie ihr Leben.}
\end{strophe}

\begin{strophe}
\vers{Brauner Staub, Du qualvoll Elend,}
\vers{Du, der genießt den Anblick}
\vers{unzähliger Seelen.}
\end{strophe}

\begin{strophe}
\vers{Dein Spuk ist zu hören}
\vers{in Stadt sowie Land.}
\vers{Dein Gesicht ist zerfressen,}
\vers{doch liegend am Straßenrand,}
\vers{möchts niemand erblicken,}
\vers{denn verhüllt von einem gewaltigen Gewand}
\vers{voll Schmach und Verachtung}
\vers{stößt es alle nur ab}
\vers{wie ein schwarzes Loch der Schande.}
\end{strophe}

\begin{strophe}
\vers{Brauner Staub, Du qualvoll Elend,}
\vers{Du, der genießt den Widerspruch}
\vers{in unzähligen Fällen.}
\end{strophe}

\begin{strophe}
\vers{Warum scheinst Du wie die Sonne?}
\vers{und bist düster wie ein Schatten?}
\vers{Wo Du auch ankommst lässt Du es blühen,}
\vers{doch hinterlässt das Gebiet zerstört, voll Ruinen.}
\end{strophe}

\begin{strophe}
\vers{Ein bisschen Staub, wie Du,}
\vers{so mächtig und teuer,}
\vers{Du bist wahrlich Mysterium,}
\vers{Du qualvoll Elend.}
\end{strophe}

\enlargethispage{1\baselineskip}

13.02.2026, ergänzt und überarbeitet am 29.06.2026

%Christoph Koloska
